\begin{result}{E.2.29}{An integral IA lift that moves actual boundary homology}{Proved one-cover counterexample}
Let \(U=\mathbb P^1_{\overline{\Q}}\setminus\{0,1,\infty\}\), fix
\(F=\langle X,Y\rangle=\pi_1^{\mathrm{top}}(U(\mathbb C))\),
and specify the quotient
\[
 \pi:\widehat F\longrightarrow H=A_5,\qquad
 X\longmapsto x=(123),\quad Y\longmapsto y=(12345).
 \tag{E.2.29a}
\]
Products of permutations act as compositions of functions. Let
\(Y_H\to U\) be its actual connected Galois cover, let \(C_H\)
be its smooth compactification, and put
\[
 \begin{aligned}
 L&=H_1(Y_H(\mathbb C),\Z),\\
 P&=\langle\text{oriented boundary classes}\rangle\subset L,\\
 W&=H_1(C_H(\mathbb C),\Z)=L/P .
 \end{aligned}
\]
There is an actual continuous automorphism
\(\alpha\in\operatorname{Aut}(\widehat F)\) with
\(\pi\alpha=\pi\) and \(\alpha^{\mathrm{ab}}=1\) such that
its action on
\(\ker(\pi)^{\mathrm{ab}}=\widehat L\) is the completion of
an integral \(H\)-equivariant automorphism \(\gamma\) of \(L\),
but \(\gamma(P)\not\subset P\). In particular this action
does not induce an endomorphism of the fixed proper-homology
quotient \(W\).

The construction has
\[
 g(C_H)=9,\quad \#(C_H\setminus Y_H)=44,\quad
 \operatorname{rank}L=61,\quad
 \operatorname{rank}P=43,\quad
 \operatorname{rank}W=18 .
 \tag{E.2.29b}
\]
It supplies an integral matrix \(N\), an actual boundary class
\(b\in P\), and an integral functional \(q:L\to\Z\) with
\[
 N^2=0,\quad \gamma=I+60N,\quad
 q(P)=0,\quad qNb=19440,\quad q\gamma b=1166400 .
 \tag{E.2.29c}
\]
The same failure holds for every positive power of \(\alpha\).
This is a statement about one specified cover. It gives no
simultaneous transporter of the proper integral lattices in
E.2.22; its failure to preserve the actual cusp module also
excludes it from the Galois image.
\end{result}
\begin{proof}
We construct the lattice maps from the covering graph, give
the extension correction explicitly, and then apply the
specified profinite lifting theorem.

\paragraph{1. The cover and the exact boundary quotient.}
The two permutations generate \(A_5\), and
\(z=(xy)^{-1}\) has order~5. The finite-cover correspondence
gives \(Y_H\); the topological and étale cover correspond under
the fixed complex embedding. The boundary has \(20,12,12\)
points over \(0,1,\infty\), respectively. Riemann--Hurwitz gives
\[
 2g(C_H)-2
 =60\left(-2+\frac23+\frac45+\frac45\right)=16.
\]
The oriented exact sequence of the compactification pair is
\[
 0\longrightarrow\Z
 \xrightarrow{\,1\mapsto\sum_{d\in D}d\,}
 \Z[D]\longrightarrow L\longrightarrow W
 \longrightarrow0,\qquad D=C_H\setminus Y_H.
 \tag{E.2.29d}
\]
It is \(H\)-equivariant, and \(P\) is saturated since \(W\)
is free. These maps prove all ranks in (E.2.29b).
This is the \(A_5\) monodromy of
\cite[Example~5.3, p.~13]{AuditBleherChinburgLubotzky2021v3};
we give the integral maps realizing the shared-constituent
phenomenon of their Theorem~5.1.

\paragraph{2. A specified integral graph basis.}
The lifted two-circle spine has vertices \(h\in H\) and
oriented edges \((h,t):h\to ht\), for \(t=x,y\). Its actual
chain sequence is
\[
 0\longrightarrow L\longrightarrow\Z[H]^2
 \xrightarrow{\ d\ }\Z[H]\xrightarrow{\mathrm{aug}}\Z
 \longrightarrow0,\qquad d(h,t)=ht-h.
 \tag{E.2.29e}
\]
The deck action is left multiplication, which commutes
with the right multiplication defining the edges.
Enumerate \(H\) by breadth-first search from \(1\), examining
\(x\) before \(y\), and retain the first edge reaching each
new vertex. This gives a tree with \(59\) edges. Order all
\(120\) edges by the resulting vertex order, taking the
\(x\)-edge before the \(y\)-edge. Let \(p_h\) be the tree path
from \(1\) to \(h\). For the \(61\) non-tree edges, in this
order, use the cycles
\[
 c_{(h,t)}=p_h+(h,t)-p_{ht}.
 \tag{E.2.29f}
\]
Their \(120\)-by-\(61\) matrix \(C_{\mathrm{cyc}}\) restricts
to the identity on non-tree rows; that restriction is an
integral inverse on the cycle module. The actual left
permutation of edges therefore gives the integral matrices
\(\rho(h)\) on \(L\), characterized by
\(C_{\mathrm{cyc}}\rho(h)=hC_{\mathrm{cyc}}\).

For each left coset \(h\langle x\rangle\), lift \(X^3\)
from \(h\); for \(h\langle y\rangle\), lift \(Y^5\);
for \(h\langle z\rangle\), lift
\((Y^{-1}X^{-1})^5\). Reversed edges have negative
coefficients. Take the first unused representative in the
same group order. This gives the actual cusp-cycle matrix
\[
 B:\Z^{44}\longrightarrow L=\Z^{61}.
 \tag{E.2.29g}
\]
Its columns sum to zero, have rank~43, and have image \(P\)
by (E.2.29d). Let \(b=B_{\!,0}\), with zero-based indexing:
it is the \(X^3\)-cusp at the identity.

\paragraph{3. The literal boundary-to-standard maps.}
Let
\[
 V=\{(a_1,\ldots,a_5)\in\Q^5:\textstyle\sum_i a_i=0\}
\]
with basis \(f_i=e_i-e_5\), \(1\leq i\leq4\), and let
\(v(h)\) be its integral permutation matrices. The
two-transitive five-letter action shows that \(V\) is
absolutely irreducible: the endomorphism algebra of the
permutation representation has dimension two, and its
trivial constituent occurs once. Sequence (E.2.29e) gives
\([L_\Q]=[\Q[H]]+[\Q]\). Sequence (E.2.29d) gives
\[
 [P_\Q]=[\Q[H/\langle x\rangle]]
       +[\Q[H/\langle y\rangle]]
       +[\Q[H/\langle z\rangle]]-[\Q].
 \tag{E.2.29h}
\]
The dimensions \( \dim V^{\langle x\rangle}=2\) and
\(\dim V^{\langle y\rangle}=\dim V^{\langle z\rangle}=0\)
show that \(V\) has multiplicities \(4,2,2\) in
\(L_\Q,P_\Q,W_\Q\), respectively. The certificate also
checks these multiplicities against the actual deck matrices.

The vector \(u=(1,1,1,-3,0)\) is \(x\)-fixed. Define
the integral matrix \(F_D:\Z[D]\to V\) by
\[
 F_D(h\langle x\rangle)=v(h)u,\qquad
 F_D(h\langle y\rangle)=F_D(h\langle z\rangle)=0 .
 \tag{E.2.29i}
\]
It is well-defined and equivariant. Its sum on the boundary
points is an invariant vector of \(V\), hence zero. It
therefore factors through the actual lattice \(P\), and
sends \(b\) to \(u\).

Here is a full finite prescription for the remaining
matrices. Let \(B_0\) be the first \(43\) columns of \(B\).
The following zero-based row indices give a unimodular
\(43\)-by-\(43\) minor \(B_{\mathcal R}\):
\[
 \begin{aligned}
 \mathcal R={}&(0,1,\ldots,29,\,
 31,32,33,34,35,36,38,43,44,50,51,54,58).
 \end{aligned}
 \tag{E.2.29j}
\]
Let \(J:\Z^{61}\to\Z^{43}\) equal
\(B_{\mathcal R}^{-1}\) on these coordinates and zero on
the others. Thus \(JB_0=I\). Let \(F_0\) be the corresponding
\(43\) columns of \(F_D\), and set
\[
 \begin{aligned}
 P_0&=B_0J:L\longrightarrow P,&
 T_0&=F_0J:L\longrightarrow V,\\
 E&=\sum_{h\in H}\rho(h)P_0\rho(h)^{-1},&
 Q_0&=60I-E .
 \end{aligned}
 \tag{E.2.29k}
\]
The finite certificate verifies \(P_0B=B\) and \(T_0B=F_D\).
Every summand defining \(E\) has image in \(P\) and is
the identity on \(P\). Consequently
\[
 E^2=60E,\qquad EB=60B,\qquad EQ_0=Q_0E=0;
 \tag{E.2.29l}
\]
moreover \(E\) and \(Q_0\) are \(H\)-equivariant.

Let \(A_0:V\to L\) send \(f_1\) to the second non-tree
cycle, of index~1, and send the other basis vectors to zero.
Set
\[
 \begin{aligned}
 A_H&=\sum_{h\in H}\rho(h)A_0v(h)^{-1}:V\longrightarrow L,\\
 S&=Q_0A_H:V\longrightarrow L,\qquad
 T=T_0E:L\longrightarrow V,\qquad
 N=\frac1{10}ST:L\longrightarrow L .
 \end{aligned}
 \tag{E.2.29m}
\]
The maps \(S,T\) are equivariant. For \(T\), this uses
that \(E(L)\subset P\) and \(T_0|_P=F_D|_P\).
Also \(TS=T_0EQ_0A_H=0\), so \(N^2=0\).
The denominator is certified integrally: the entries of
\(ST\) have greatest common divisor~10. Thus \(N\) is
an actual integral \(61\)-by-\(61\) matrix of rank~4,
with largest absolute entry \(3156\).

Let \(q\) be the first row of \(Q_0\). The exact
integer evaluations are
\[
 qB=0,\qquad (Nb)_0=324,\qquad qNb=19440 .
 \tag{E.2.29n}
\]
All matrices in (E.2.29f)--(E.2.29n), including the full
\(61\)-by-\(61\) matrix \(N\), are retained in
\texttt{code/build\_A5\_boundary\_}\allowbreak
\texttt{nilpotent.json}. The self-contained program with
the same stem and suffix \texttt{.py} reconstructs them
from the two specified permutations. It uses exact
integer and rational arithmetic; every potentially large
fixed-width matrix product has a proved numerical bound
before multiplication, and an independent arbitrary-precision
recheck verifies the decisive identities.

It follows that
\(\gamma=I+60N\) and \(\gamma^{-1}=I-60N\)
are mutually inverse integral \(H\)-maps. Since \(q\)
factors through \(W\), (E.2.29n) gives (E.2.29c).
Any quotient square \(W\to W\) would send the class of
\(b\), which is zero, to the nonzero class of \(\gamma b\).
This proves the asserted failure of descent through the
actual quotient map.

\paragraph{4. The extension cocycle and its explicit correction.}
Put \(K=\ker(F\to H)\). The graph identifies
\(K/[K,K]=L\). For the specified tree word \(w_h\),
the section \(h\mapsto w_h\) has factor set
\[
 c(h,k)=[w_hw_kw_{hk}^{-1}]
       =[p_h+h p_k-p_{hk}]\in L.
 \tag{E.2.29o}
\]
The last equality is the actual chain-to-abelianization
map. Non-tree restriction computes its coordinates; the
two unshifted tree paths contribute zero there.
This gives the discrete extension
\(1\to L\to F/[K,K]\to H\to1\).

The finite-index subgroup \(K\) has its full induced
profinite topology: the core in \(F\) of a finite-index
subgroup of \(K\) remains finite index. Thus
\(R=\ker(\widehat F\to H)=\widehat K\), and its maximal
abelian profinite quotient is \(\widehat L\).
Throughout this paragraph \([R,R]\) denotes the
\emph{closed} commutator subgroup. The same factor set
represents the extension
\[
 1\longrightarrow\widehat L\longrightarrow
 \widehat F/[R,R]\longrightarrow H\longrightarrow1,
 \quad \beta\in H^2(H,\widehat L).
 \tag{E.2.29p}
\]
Define
\[
 b_0(h)=\sum_{k\in H}c(h,k),\qquad \eta(h)=Nb_0(h).
 \tag{E.2.29q}
\]
Summing the cocycle identity
\(c(h,k)+c(hk,l)=h c(k,l)+c(h,kl)\) over \(l\)
gives the exact sign and identities
\[
 \delta b_0(h,k)=h b_0(k)-b_0(hk)+b_0(h)=60c(h,k),
 \qquad
 \delta\eta=(\gamma-I)c.
 \tag{E.2.29r}
\]
The program checks (E.2.29r) on all \(3600\) pairs
and retains the two cochains. In section coordinates
the group law and the actual automorphism are
\[
 \begin{aligned}
 (m,h)(n,k)&=(m+h n+c(h,k),hk),\\
 \widetilde\gamma(m,h)&=(\gamma m+\eta(h),h),\\
 \widetilde\gamma^{-1}(m,h)&=((I-60N)m-\eta(h),h).
 \end{aligned}
 \tag{E.2.29s}
\]
The multiplication check is exactly (E.2.29r).
The inverse uses \(N\eta=N^2b_0=0\).
These formulas are continuous after completing \(L\).
In particular \(\widehat\gamma\) fixes \(\beta\), with
an explicit correction cochain, and
\(\widetilde\gamma\) induces the identity on \(H\).

\paragraph{5. The actual lift and its base abelianization.}
Bleher--Chinburg--Lubotzky
\cite[Theorem~2.1, p.~4; Corollary~2.4 and its proof,
p.~5]{AuditBleherChinburgLubotzky2021v3}
gives
\[
 \{\alpha\in\operatorname{Aut}(\widehat F):\pi\alpha=\pi\}
 \longrightarrow
 \operatorname{Aut}_{H,\beta}(\widehat L)
 \quad\text{surjectively}.
 \tag{E.2.29t}
\]
More precisely, their Gaschütz corollary lifts the
specified \(\widetilde\gamma\) on
\(\widehat F/[R,R]\) to an actual automorphism
\(\alpha\) of \(\widehat F\). Its quotient diagram
gives \(\pi\alpha=\pi\), and its action on
\(R^{\mathrm{ab}}\) is exactly \(\widehat\gamma\).
The theorem supplies the profinite lift; no infinite
profinite generator words are claimed to be computed.

Let \(A_{\mathrm{exp}}:L\to\Z^2\) sum respectively
the \(X\)-edge and \(Y\)-edge coefficients. This is the
actual map induced by \(K\hookrightarrow F\).
It is \(H\)-invariant. Since the image of \(N\) factors
through the nontrivial irreducible \(V\),
\(A_{\mathrm{exp}}N=0\); the matrix identity is also
checked integrally. The cycle columns of indices
\(1,6\) give
\[
 (A_{\mathrm{exp}})_{\{1,6\}}
 =\begin{pmatrix}2&1\\-1&-1\end{pmatrix},
 \qquad\det=-1.
 \tag{E.2.29u}
\]
Their literal based words are
\[
 r=XYYX(YYY)^{-1},\qquad
 s=YXXYX(XYXYY)^{-1}.
 \tag{E.2.29v}
\]
Both evaluate to \(1\) in the specified \(A_5\), as
the certificate verifies. The words \(rs^{-1}\)
and \(rs^{-2}\) have exponent vectors \((1,0)\)
and \((0,1)\). Hence \(A_{\mathrm{exp}}\) is onto
integrally and its completion is onto. Naturality for
the actual inclusion \(R\hookrightarrow\widehat F\) gives
\[
 \alpha^{\mathrm{ab}}\widehat A_{\mathrm{exp}}
 =\widehat A_{\mathrm{exp}}\widehat\gamma
 =\widehat A_{\mathrm{exp}},
 \quad\text{so}\quad \alpha^{\mathrm{ab}}=1 .
 \tag{E.2.29w}
\]
This proves the IA strengthening for every lift with
the specified relation-module action.

Finally, \(\gamma^m=I+60mN\) gives
\(q\gamma^m b=1166400m\ne0\) for every positive integer
\(m\). Every finite-index subgroup of the quotient-fixing
automorphism group contains such a power of \(\alpha\);
thus it cannot preserve this peripheral lattice.
A Galois automorphism fixing the specified cover quotient
permutes its actual cusp inertia subgroups and acts on
their generators by units. Its abelianized map therefore
preserves \(\widehat P\). The nonzero value in
(E.2.29c), also nonzero in \(\widehat\Z\), proves
directly that the constructed lift cannot have that
Galois property. The example consequently tests the
one-cover profinite-lifting inference without asserting
the all-cover hypothesis required in E.2.22 or E.2.28.
\end{proof}
